\chapter{Initial and Boundary Value Problems}

% -------------------------------------------------------
\section{Initial Value Problem}
% -------------------------------------------------------

\begin{definition}[Initial Value Problem]
An \textbf{initial value problem} (IVP) is a differential equation together with conditions that specify the value of the solution (and possibly its derivatives) at a \emph{single point}.
\end{definition}

\begin{example}
\[
  \frac{dy}{dx} = f(x,y), \qquad y(x_0) = y_0
\]
The condition $y(0) = 1$ is a typical initial condition.
\end{example}

% -------------------------------------------------------
\section{Boundary Value Problem}
% -------------------------------------------------------

\begin{definition}[Boundary Value Problem]
A \textbf{boundary value problem} (BVP) is a differential equation together with conditions specified at \emph{two or more distinct points}.
\end{definition}

\begin{example}
\[
  y'' + y = 0, \qquad y(0) = 0,\quad y(\pi) = 1
\]
The conditions are prescribed at both endpoints of the interval $[0,\pi]$.
\end{example}

\begin{remark}
IVPs arise naturally in time-evolution problems (all conditions given at $t = 0$), while BVPs arise in spatial problems (conditions given at the boundaries of a region).
\end{remark}
